% FAC-EN-U0001
% French source: fac_body.tex lines 1--115.
% Authority: PDF physical pp. 2--5; printed pp. 197--200 (through section 2).

\begin{center}
\textsc{Annals of Mathematics}

Vol. 61, No. 2, March 1955

\textit{Printed in U.S.A.}

\vspace{1.2em}
{\Large\bfseries COHERENT ALGEBRAIC SHEAVES}

\vspace{1em}
\textsc{By Jean-Pierre Serre}

(Received October 8, 1954)
\end{center}

\phantomsection\label{fac:introduction}
\section*{Introduction}
\addcontentsline{toc}{section}{Introduction}

It is well known that cohomological methods, and sheaf theory in particular, play an increasingly important role, not only in the theory of functions of several complex variables (cf. [5]), but also in classical algebraic geometry (it will suffice to cite the recent work of Kodaira and Spencer on the Riemann--Roch theorem). The algebraic character of these methods suggested that they could also be applied to abstract algebraic geometry; the purpose of the present paper is to show that this is indeed the case.

The contents of the various chapters are as follows:

\hyperref[fac:I]{Chapter I} is devoted to the general theory of sheaves. It contains proofs of the results from that theory which are used in the other two chapters. The various algebraic operations that may be performed on sheaves are described in \hyperref[fac:I:1]{\S1}; we have followed Cartan's account ([2], [5]) fairly closely. \hyperref[fac:I:2]{Section 2} studies coherent sheaves of modules; these sheaves generalize coherent analytic sheaves (cf. [3], [5]) and enjoy entirely analogous properties. In \hyperref[fac:I:3]{\S3} the cohomology groups of a space $X$ with coefficients in a sheaf $\mathfrak{F}$ are defined. In the later applications, $X$ is an algebraic variety equipped with the Zariski topology; it is therefore not a Hausdorff topological space, and the methods used by Leray [10] or Cartan [3] (based on ``partitions of unity'' or ``fine'' sheaves) do not apply to it. We have therefore had to return to the \v{C}ech procedure and define the cohomology groups $H^q(X,\mathfrak{F})$ by passage to the direct limit over increasingly fine open coverings. A further difficulty, related to the non-Hausdorffness of $X$, arises in the ``cohomology exact sequence'' (cf. nos.~\hyperref[fac:I:3:no:24]{24} and \hyperref[fac:I:3:no:25]{25}): we have been able to establish this exact sequence only in certain special cases, which are nevertheless sufficient for the applications we had in view (cf. nos.~\hyperref[fac:I:3:no:24]{24} and \hyperref[fac:II:3:no:47]{47}).

\hyperref[fac:II]{Chapter II} begins with a definition of algebraic varieties analogous to Weil's ([17], Chap.~VII), but also encompassing reducible varieties. We note in this connection that, contrary to Weil's usage, we do not reserve the term ``variety'' for irreducible varieties. We define the structure of an algebraic variety by specifying a topology (the Zariski topology) and a subsheaf of the sheaf of germs of functions (the sheaf of local rings). A coherent algebraic sheaf on an algebraic variety $V$ is simply a coherent sheaf of $\mathcal{O}_V$-modules, where $\mathcal{O}_V$ denotes the sheaf of local rings of $V$; we give various examples in \hyperref[fac:II:2]{\S2}. \hyperref[fac:II:3]{Section 3} is devoted to affine varieties. The results obtained are entirely analogous to those concerning Stein varieties (cf. [3], [5]): if $\mathfrak{F}$ is a coherent algebraic sheaf on the affine variety $V$, then $H^q(V,\mathfrak{F})=0$ for every $q>0$, and $\mathfrak{F}_x$ is generated by $H^0(V,\mathfrak{F})$ for every $x\in V$. Moreover (\hyperref[fac:II:4]{\S4}), $\mathfrak{F}$ is uniquely determined by $H^0(V,\mathfrak{F})$, regarded as a module over the coordinate ring of $V$.

\hyperref[fac:III]{Chapter III}, concerning projective varieties, contains the essential results of this paper. We begin by establishing a correspondence between coherent algebraic sheaves $\mathfrak{F}$ on projective space $X=\mathbf{P}_r(K)$ and graded $S$-modules satisfying condition (TF) of no.~\hyperref[fac:III:2:no:56]{56} (where $S$ denotes the polynomial algebra $K[t_0,\ldots,t_r]$). This correspondence is bijective if two $S$-modules which differ only in their homogeneous components of sufficiently low degree are identified (for the precise statements, see nos.~\hyperref[fac:III:2:no:57]{57}, \hyperref[fac:III:2:no:59]{59}, and \hyperref[fac:III:3:no:65]{65}). From that point on, every question concerning $\mathfrak{F}$ can be translated into a question concerning the associated $S$-module $M$. Thus, in \hyperref[fac:III:3]{\S3} we give a procedure for determining the $H^q(X,\mathfrak{F})$ algebraically from $M$. This enables us, in particular, to study the properties of $H^q(X,\mathfrak{F}(n))$ as $n\to+\infty$ (for the definition of $\mathfrak{F}(n)$, see no.~\hyperref[fac:III:2:no:54]{54}); the resulting statements are given in nos.~\hyperref[fac:III:3:no:65]{65} and \hyperref[fac:III:3:no:66]{66}. In \hyperref[fac:III:4]{\S4} we relate the groups $H^q(X,\mathfrak{F})$ to the functors $\operatorname{Ext}^q_S$ introduced by Cartan and Eilenberg [6]. This allows us, in \hyperref[fac:III:5]{\S5}, to study the behavior of $H^q(X,\mathfrak{F}(n))$ as $n\to-\infty$ and to give a homological characterization of varieties that are ``$k$-fold of the first kind.'' \hyperref[fac:III:6]{Section 6} presents some elementary properties of the Euler--Poincar\'e characteristic of a projective variety with coefficients in a coherent algebraic sheaf.

Elsewhere we shall show how the general results of the present paper can be applied to various particular problems, notably to extend the ``duality theorem'' of [15] to the abstract case, as well as part of Kodaira and Spencer's results on the Riemann--Roch theorem; in these applications the theorems of nos.~\hyperref[fac:III:3:no:66]{66}, \hyperref[fac:III:5:no:75]{75}, and \hyperref[fac:III:5:no:76]{76} play an essential role. We shall also show that, when the base field is the field of complex numbers, the theory of coherent algebraic sheaves is essentially identical to that of coherent analytic sheaves (cf. [4]).

\begin{center}
\textsc{Table of Contents}
\end{center}

\noindent\hyperref[fac:I]{\textsc{Chapter I. Sheaves}}

\noindent\hyperref[fac:I:1]{\S1}. Operations on sheaves\dotfill 199

\noindent\hyperref[fac:I:2]{\S2}. Coherent sheaves of modules\dotfill 207

\noindent\hyperref[fac:I:3]{\S3}. Cohomology of a space with coefficients in a sheaf\dotfill 212

\noindent\hyperref[fac:I:4]{\S4}. Comparison of the cohomology groups of different coverings\dotfill 219

\medskip
\noindent\hyperref[fac:II]{\textsc{Chapter II. Algebraic varieties---Coherent algebraic sheaves on affine varieties}}

\noindent\hyperref[fac:II:1]{\S1}. Algebraic varieties\dotfill 223

\noindent\hyperref[fac:II:2]{\S2}. Coherent algebraic sheaves\dotfill 230

\noindent\hyperref[fac:II:3]{\S3}. Coherent algebraic sheaves on affine varieties\dotfill 233

\noindent\hyperref[fac:II:4]{\S4}. Correspondence between finitely generated modules and coherent algebraic sheaves\dotfill 240

\medskip
\noindent\hyperref[fac:III]{\textsc{Chapter III. Coherent algebraic sheaves on projective varieties}}

\noindent\hyperref[fac:III:1]{\S1}. Projective varieties\dotfill 243

\noindent\hyperref[fac:III:2]{\S2}. Graded modules and coherent algebraic sheaves on projective space\dotfill 246

\noindent\hyperref[fac:III:3]{\S3}. Cohomology of projective space with coefficients in a coherent algebraic sheaf\dotfill 253

\noindent\hyperref[fac:III:4]{\S4}. Relations with the functors $\operatorname{Ext}^q_S$\dotfill 261

\noindent\hyperref[fac:III:5]{\S5}. Applications to coherent algebraic sheaves\dotfill 267

\noindent\hyperref[fac:III:6]{\S6}. Characteristic function and arithmetic genus\dotfill 274

\begin{center}
\phantomsection\label{fac:I}
{\Large\textsc{Chapter I. Sheaves}}
\end{center}

\phantomsection\label{fac:I:1}
\section*{\S1. Operations on sheaves}
\addcontentsline{toc}{section}{Chapter I, \S1. Operations on sheaves}

\phantomsection\label{fac:I:1:no:1}
\textbf{1. Definition of a sheaf.} Let $X$ be a topological space. A \textit{sheaf of abelian groups} on $X$ (or simply a \textit{sheaf}) consists of:

\begin{enumerate}[label=(\alph*),leftmargin=2.5em]
\item\label{fac:I:1:no:1:data:item:a} \textit{A function $x\mapsto\mathfrak{F}_x$ assigning to each $x\in X$ an abelian group $\mathfrak{F}_x$,}
\item\label{fac:I:1:no:1:data:item:b} \textit{A topology on the set $\mathfrak{F}$, the disjoint union of the sets $\mathfrak{F}_x$.}
\end{enumerate}

If $f$ is an element of $\mathfrak{F}_x$, we put $\pi(f)=x$; the map $\pi$ is called the \textit{projection} of $\mathfrak{F}$ onto $X$. The subset of $\mathfrak{F}\times\mathfrak{F}$ formed by the pairs $(f,g)$ such that $\pi(f)=\pi(g)$ will be denoted by $\mathfrak{F}+\mathfrak{F}$.

With these definitions in place, we can state the two axioms to which the data \hyperref[fac:I:1:no:1:data:item:a]{(a)} and \hyperref[fac:I:1:no:1:data:item:b]{(b)} are subject:

\begin{enumerate}[label=(\Roman*),leftmargin=3em]
\item\label{fac:I:1:no:1:axiom:I} \textit{For every $f\in\mathfrak{F}$ there exist a neighborhood $V$ of $f$ and a neighborhood $U$ of $\pi(f)$ such that the restriction of $\pi$ to $V$ is a homeomorphism of $V$ onto $U$.}

(In other words, $\pi$ must be a local homeomorphism.)

\item\label{fac:I:1:no:1:axiom:II} \textit{The map $f\mapsto-f$ is continuous from $\mathfrak{F}$ to $\mathfrak{F}$, and the map $(f,g)\mapsto f+g$ is continuous from $\mathfrak{F}+\mathfrak{F}$ to $\mathfrak{F}$.}
\end{enumerate}

Observe that, even if $X$ is Hausdorff (which we have not assumed), $\mathfrak{F}$ need not be Hausdorff, as is shown by the example of the sheaf of germs of functions (cf. no.~\hyperref[fac:I:1:no:3]{3}).

\textsc{Example} of a sheaf. Let $G$ be an abelian group and put $\mathfrak{F}_x=G$ for every $x\in X$. The set $\mathfrak{F}$ may be identified with the product $X\times G$; if it is given the product topology of the topology on $X$ and the discrete topology on $G$, one obtains a sheaf, called the \textit{constant sheaf} isomorphic to $G$, and often identified with $G$.

\medskip
\phantomsection\label{fac:I:1:no:2}
\textbf{2. Sections of a sheaf.} Let $\mathfrak{F}$ be a sheaf on the space $X$, and let $U$ be a subset of $X$. A \textit{section} of $\mathfrak{F}$ over $U$ is a continuous map $s:U\to\mathfrak{F}$ such that $\pi\circ s$ is the identity map of $U$. Thus $s(x)\in\mathfrak{F}_x$ for every $x\in U$. The set of sections of $\mathfrak{F}$ over $U$ will be denoted by $\Gamma(U,\mathfrak{F})$; axiom \hyperref[fac:I:1:no:1:axiom:II]{(II)} implies that $\Gamma(U,\mathfrak{F})$ is an abelian group. If $U\subset V$ and $s$ is a section over $V$, the restriction of $s$ to $U$ is a section over $U$; this gives a homomorphism $\rho_U^V:\Gamma(V,\mathfrak{F})\to\Gamma(U,\mathfrak{F})$.

If $U$ is open in $X$, then $s(U)$ is open in $\mathfrak{F}$; and, as $U$ ranges over a basis of open subsets of $X$, the sets $s(U)$ range over a basis of open subsets of $\mathfrak{F}$. This is simply another way of expressing axiom \hyperref[fac:I:1:no:1:axiom:I]{(I)}.

Let us also note a consequence of axiom \hyperref[fac:I:1:no:1:axiom:I]{(I)}: for every $f\in\mathfrak{F}_x$, there exists a section $s$ over a neighborhood of $x$ such that $s(x)=f$, and any two sections with this property coincide in a neighborhood of $x$. In other words, $\mathfrak{F}_x$ is the direct limit of the $\Gamma(U,\mathfrak{F})$ over the directed set of neighborhoods of $x$.
